\paragraph{Cooperation and Mutual Help as the Central Driver}
Across all experimental settings, our simulations consistently show that \emph{cooperation and mutual help are the primary drivers of evolutionary survival} (Table~\ref{tab:replicate_outcomes_main}). Universal and reciprocal agents---the most cooperative types---exhibit the most stable survival rates across all conditions. Kin-focused agents remain viable when resources and communication bandwidth suffice for them to reach self-sustaining family size. Selfish agents are strongly disfavoured in every setting. These results are broadly consistent with prior theoretical predictions~\citep{hamilton1964genetical,trivers1971evolution,nowak1998dynamics,Curry2019MAC}, and support the expanding circle model~\citep{singer1981expanding}: broader moral circles generally produce superior evolutionary outcomes.

\paragraph{What the Paradigm Unlocks}
The two mechanisms isolated in \S\ref{sec:mechanisms} are the direct payoff of relaxing both simplifications of prior work (Fig.~\ref{fig:paradigm}): self-purging requires agents who recognise same-type competitors; judgment cost requires agents who take time to form beliefs about one another and can misjudge. Action-level resolution further lets us trace \emph{why} a type fails by examining individual agents' reasoning (supplement \S\ref{sec:case_studies}). These patterns also echo anthropological evidence that kin-focused cooperation~\citep{holden2003matriliny, mattison2011evolutionary} preceded the expansion of moral circles under environmental pressure~\citep{singer1981expanding}.

\paragraph{Conclusions}
We propose a new research paradigm that brings cognitive realism to the study of moral evolution through LLM-based agent simulation. Our experiments show that cooperation and mutual help are the central driver of evolutionary survival, with universal and reciprocal morality exhibiting the most stable outcomes. Two mechanisms emerge from cognition---self-purging of selfish agents and the cost of moral judgment---that decisively shape evolutionary outcomes and are inaccessible to traditional approaches. We hope this paradigm can inspire broader and deeper investigations into moral evolution and other complex social phenomena that require cognitively realistic modeling.
